En A la velocidad satisface una ecuación de primer orden y vale
\(y'=(1-e^{-2t})/2\): la curva parte horizontal, crece, es cóncava hacia arriba y su pendiente tiende a \(1/2\). En B la respuesta oscila entre \(0\) y \(1\). Ambas parten con aceleración uno, pero ese comienzo común no determina el comportamiento posterior.
